\begin{titlepage}
\centering
{\Large Notas didácticas de entrevista a fondo\par}
\vspace{1.0cm}
{\huge\bfseries Manus: la última entrevista, la travesía de los Agent y los límites de la complejidad\par}
\vspace{0.5cm}
{\Large Zhang Xiaojun conversa con Ji Yichao (季逸超) Peak\par}
\vspace{0.35cm}
{\large Elaboradas a partir del video público de Zhang Xiaojun Podcast, la transcripción hecha con GPU, la estructura de capítulos y diagramas conceptuales propios\par}
\vspace{0.3cm}
{\large 2025-12-30\par}
\vspace{0.85cm}
\includegraphics[width=0.88\textwidth,height=0.42\textheight,keepaspectratio]{cover.jpg}\par
\vfill
\begin{tcolorbox}[width=0.92\textwidth, colback=black!2!white, colframe=black!55, sharp corners]
\textbf{Título del video}: 128. La última entrevista antes de que Manus decidiera venderse: ¡ah, esta fantástica travesía de 2025…!\par
\textbf{Canal del video}: Zhang Xiaojun Podcast\par
\textbf{Duración del video}: 03:31:17\par
\textbf{Enlace del video}: \href{\videourl}{\nolinkurl{\videourl}}\par
\textbf{Nota sobre los materiales}: YouTube no ofrece subtítulos disponibles; el texto principal se elaboró a partir de la transcripción con faster-whisper large-v3 en CUDA, los capítulos del video, la portada y figuras didácticas propias. Las tomas fijas de la entrevista no se repiten en el texto principal.\par
\textbf{Página del proyecto}: \href{\repourl}{\nolinkurl{\repourl}}\par
\end{tcolorbox}
\end{titlepage}

\tableofcontents
\newpage
